\documentclass{article}
\begin{document}
\usepackage[utf8]{inputenc}
\usepackage[T1]{fontenc}
\usepackage{amsmath}
\usepackage{amssymb}
\usepackage{amsthm}
\usepackage{geometry}
\usepackage{graphicx}
\usepackage{hyperref}
\usepackage{natbib}
\geometry{margin=1in}

\title{Mind at Work: Detecting Agency through Assembly and Action}
\author{Flyxion}
\date{September 10, 2025}

\begin{document}

\maketitle

\begin{abstract}
This monograph presents a comprehensive framework for detecting mind-like agency by integrating the Effective Assembly Index for Mind Recognition (EAIMR) with the praxicon concept from cognitive science. EAIMR quantifies the improbability of a system’s structure arising from unguided processes, adjusted for efficiency and amplified by coherence, while the praxicon provides a structured repository of action representations bridging perceptual, semantic, and motor domains. By combining these, we propose a robust method to identify agentic intelligence across neural systems, collective organisms, ecosystems, and artificial intelligence, avoiding anthropocentric biases that marginalize non-narrative minds, such as those with aphantasia, anendophasia, or ecological intelligences. The framework is formalized using information theory, graph theory, category theory, and sheaf theory, with applications in neuroscience, AI alignment, ecological cognition, and cross-cultural studies. We address intersubjectivity collapse as a structural risk, amplified by Big Tech’s attention-compute arms race and ecological fragmentation, drawing parallels to historical critiques of industrialization. The monograph contrasts this approach with competing theories, provides detailed case studies, and discusses ethical implications for recognizing diverse minds. This extended edition includes multi-scale temporal dynamics, causal intervention frameworks, compositional agency, quantitative validation studies, computational implementations, statistical methods, phenomenological bridges, evolutionary dynamics, information-theoretic refinements, AI alignment protocols, clinical applications, policy recommendations, bias mitigation strategies, computational tractability improvements, and connections to plant cognition, collective intelligence, network neuroscience, indigenous knowledge, and extended cognition.
\end{abstract}

\section{Introduction}
\label{sec:introduction}

The quest to detect and understand agency---defined as the capacity of a system to pursue goals through coordinated actions deviating from unguided physical processes---spans cognitive science, philosophy of mind, artificial intelligence (AI), and ecology. Our civilization’s bias toward equating intelligence with narrativity and imagery, dominant human cognitive modes, systematically marginalizes alternative forms of intelligence. As Michael Levin articulates, ``mind blindness'' leads us to dismiss the agency of plants, animals, and cellular collectives because they lack human-like symbols or narratives \citep{levin2022technological}. Levin's work emphasizes that intelligence is not brain-bound but present in any system capable of goal pursuit and adaptation, from cellular regeneration to collective behaviors \citep{levin2019computational}. This bias extends to cognitive minorities, such as individuals with aphantasia (inability to visualize mentally) or anendophasia (absence of inner monologue), who face exclusion in narrative-driven economies despite intact functional reasoning \citep{zeman2015, hurlburt2024}.

Big Tech’s business model---selling compute and storage while buying attention through algorithmic amplification---exacerbates this collapse. By spotlighting outlier success stories and suppressing systemic critiques, platforms create a narrative arms race that erodes trust and shared meaning \citep{zuboff2019age, wu2016attention}. Programming courses exemplify this: an oversupply of training with undersupplied job demand creates a hope market, where platforms profit by amplifying rare successes while burying failures \citep{carr2010shallows}. This dynamic parallels historical ecological critiques, from Joni Mitchell’s ``Big Yellow Taxi'' lamenting the paving of paradise for parking lots to Jim Amos’s vision of a ``Coruscant planet'' where data centers and nuclear plants displace ecosystems for a compute-driven metaverse \citep{lanham2006economics}. Roads, like digital feeds, fragment habitats---disrupting migrations, raising heat islands, and diverting water---leading to biodiversity loss, population isolation, and ecosystem degradation \citep{forman2003road}. Both create isolated patches, undermining exchange of animals or ideas, driving collapse \citep{arendt1958human}.

The Effective Assembly Index for Mind Recognition (EAIMR) and the praxicon framework offer a solution by quantifying agency through structural complexity and action representation, bypassing narrativity bias to recognize diverse intelligences. EAIMR measures the improbability of a system’s assembly, adjusted for reuse, parallelism, and degeneracy, and amplified by coherence \citep{walker2013algorithmic}. The praxicon, a cognitive ``dictionary'' of actions, bridges perceptual, semantic, and motor domains, grounding intelligence in embodied processes \citep{pastra2011praxicon, peigneux2004imaging}. Together, they detect agency in systems from neural circuits to termite mounds, forests, and AI, fostering inclusivity for non-human and atypical human minds. This monograph integrates these frameworks with information theory, graph theory, category theory, and sheaf theory, addressing applications in neuroscience, AI alignment, ecology, and cross-cultural cognition, while situating IC as a meta-crux for debates on AI futures, such as those between Yudkowsky’s pessimism and Miller’s cooperative optimism \citep{yudkowsky2008artificial, miller2006robust}.

To appreciate this integration, it is essential to understand prerequisite concepts. Information theory, pioneered by Claude Shannon, quantifies uncertainty and structure in systems through measures like entropy, which represents the average information content or disorder in a probability distribution \citep{shannon1948mathematical}. Complexity measures, such as Kolmogorov complexity, define the shortest program capable of generating a given structure, providing a basis for assessing improbability \citep{kolmogorov1965three}. Assembly theory extends these ideas to biological and technological contexts by counting the minimal steps required to construct complex objects from simple building blocks \citep{walker2013algorithmic}. The praxicon, rooted in neuropsychology, models how the brain stores and retrieves action knowledge, distinguishing between perceptual recognition and motor execution \citep{peigneux2004imaging}. By weaving these threads together, our framework addresses the challenge of detecting agency without relying on anthropocentric markers like language or self-report.

Furthermore, this work addresses ``mind blindness,'' a term from Michael Levin’s research describing the human tendency to overlook intelligence in non-human systems due to their lack of narrative or imagistic expression. This bias extends to cognitive minorities within humanity, such as individuals with aphantasia (inability to visualize mentally) or anendophasia (lack of inner speech), who process information functionally but are marginalized in narrative-dominant societies. Integrating these insights, the framework mitigates intersubjectivity collapse---the breakdown of shared understanding among diverse minds---by recognizing alternative intelligences through structural metrics like EAIMR.

\section{Historical and Conceptual Foundations}
\label{sec:foundations}

\subsection{Prerequisites: Complexity and Information in Natural Systems}
Before delving into assembly theory, consider the foundational role of complexity in distinguishing living and agentic systems from inert matter. Murray Gell-Mann’s work on effective complexity highlights how systems exhibit regularities that are neither too random nor too ordered, a balance often indicative of adaptive processes \citep{adami2012use}. Robert Hazen’s research on mineral evolution and biocomplexity introduces functional information as a metric for the emergence of life-like properties, where improbable configurations serve functional roles \citep{hazen2007functional}. These ideas set the stage for assembly theory, which quantifies the ``assembly index'' as the minimal number of operations needed to produce a structure, emphasizing reuse and hierarchy in evolved systems.

\subsection{Assembly Theory and Complexity}
Assembly theory, developed by Sara Imari Walker, Leroy Cronin, and collaborators, provides a formal measure of complexity by calculating the shortest path to assemble a molecule or structure from basic units \citep{cronin2016beyond, walker2013algorithmic}. It has applications in astrobiology and origins-of-life research, where high assembly indices correlate with biotic processes. EAIMR builds upon this by incorporating adjustments for real-world efficiencies and coherence, enabling its use as a detector for mind-like agency \citep{adami2012use}. Loopholes in assembly index application---such as reuse of substructures, parallelism, context dependence, degeneracy, and information leakage---highlight how complexity can bootstrap without full cost, relevant to understanding emergent intelligences in swarms or ecosystems.

\subsection{Praxicon: Action Representation in Cognition}
The praxicon emerges from studies of apraxia, a disorder impairing purposeful movements despite intact motor function. Prerequisite knowledge includes the distinction between ideational apraxia (conceptual deficits) and ideomotor apraxia (execution deficits), as modeled by Rothi and Heilman \citep{mantovani2010reviewing}. The praxicon is conceptualized as a dual-component system: the input praxicon for recognizing actions via visual-gestural cues, and the output praxicon for producing motor engrams \citep{peigneux2004imaging}. Extensions to AI involve grounding symbolic language in sensorimotor primitives, addressing the symbol grounding problem \citep{harnad1990symbol, pastra2011praxicon, vannuscorps2023effector}.

\subsection{Competing Approaches}
To contextualize our framework, compare it with alternatives. Integrated Information Theory (IIT) posits consciousness as integrated information $\Phi$, but it requires predefined system boundaries and struggles with distributed agency \citep{tononi2004information}. The Free Energy Principle (FEP) frames agents as minimizers of variational free energy, modeling predictive processing but lacking direct measures of structural improbability \citep{friston2010}. Enactive cognition emphasizes autopoiesis and sensorimotor contingencies, yet it often remains qualitative \citep{barandiaran2009defining}. Our EAIMR-praxicon approach complements these by providing computable metrics grounded in assembly and action representation.

\section{Mind Blindness and Intersubjectivity Collapse}
\label{sec:mind_blindness}

Mind blindness, as articulated by Michael Levin, refers to the systematic failure to recognize intelligence in non-human systems, such as plants, animals, and cellular collectives, because they do not conform to human-like symbolic or narrative modes. Levin argues that intelligence is any system’s ability to pursue goals, adapt to perturbations, and solve problems in its medium, evident in regenerating limbs or bioelectric signaling \citep{levin2019computational}. This bias extends beyond ecology to cognitive minorities within humanity, where individuals with aphantasia---the inability to voluntarily create mental images---or anendophasia---the absence of inner monologue---are marginalized despite functional processing capabilities.

\textbf{Mind Blindness Across Domains}

1. \textbf{Ecological Minds (Plants and Animals)} Bias: Intelligence is equated with language and narrative. Blindness: Plants and animals operate through signaling, adaptation, and coordination, but because they do not ``speak,'' their intelligence is dismissed as instinct. Consequence: Systematic marginalization of non-human agency, paving the way for ecological exploitation.

2. \textbf{Cognitive Minorities (Aphantasia, Anendophasia, Atypical Cognition)} Bias: Value is tied to attention performance---vivid self-narration, imagery, persuasion. Blindness: Those whose minds process functionally but without inner imagery or speech must reverse-engineer outputs, making them invisible or ``less capable'' in the attention economy. Aphantasia affects 2-4\% of the population, impairing voluntary visualization, while anendophasia involves silent inner worlds without verbal thought. Consequence: Social and economic exclusion of cognitive styles that do not match platform-optimized norms.

3. \textbf{Religious and Philosophical Systems (Christianity, Islam, Atheism, etc.)} Bias: Each tradition assumes its own narrative or epistemic mode (ritual, revelation, rational critique) is the universal medium of truth. Blindness: Values and intelligences encoded in rival traditions become illegible. What one calls wisdom, another calls heresy or superstition. Consequence: Breakdown of shared meaning across civilizations, escalating into mistrust, conflict, or war.

Unifying Thread: Across all three layers, mind blindness is the refusal to recognize the legitimacy of intelligences that do not mirror our own mode of expression. This blindness is not just a moral failure---it is a structural driver of intersubjectivity collapse, as communication, empathy, and trust fail across ecological, cognitive, and cultural domains. AI futures will magnify this bias---if we only design for intelligences that ``talk like us,'' we will fail to cooperate with (or even recognize) those that do not.

\subsection{Termite--Neuron--Forest Triad: Analogies for Detecting Emergent Intelligence}

In the context of intersubjectivity collapse---where diverse minds fail to communicate or recognize one another due to incompatible cognitive modes---the ``Termite--Neuron--Forest Triad'' serves as a formal analogy to highlight how emergent intelligence arises from swarms of simpler agents. This triad underscores the structural bias of mind blindness, wherein intelligences lacking narrative or imagistic expression are dismissed, exacerbating collapse by excluding valid forms of agency. By applying an assembly-index lens, we can quantify the improbability of these emergent structures, revealing mind-like processes that operate through coordination rather than human-centric symbolism.

1. \textbf{Neurons as ``Just Insects'': The Swarm of Consciousness} Neurons, individually, are simple electrochemical agents following local rules of excitation and inhibition. However, their collective swarming---through synaptic connections, feedback loops, and hierarchical organization---assembles a brain capable of consciousness and goal-directed behavior. The assembly index here is extraordinarily high: the precise mesoscale architecture (e.g., cortical columns) is improbable under mindless diffusion, implying directed processes like plasticity and learning. Mind blindness dismisses this as ``mere biology'' until the emergent whole is recognized, mirroring how we overlook atypical human cognition (e.g., aphantasia) that lacks internal imagery.

2. \textbf{Termites as ``Just Insects'': The Swarm of Architecture} Termites operate as decentralized agents, responding to local pheromones and environmental cues. Yet, their swarming constructs mounds with advanced features: ventilation topology for thermoregulation, chambers for farming fungi, and resilient structures against floods. The mound’s assembly index exceeds that of random soil accumulations (e.g., dunes), reflecting collective intelligence in distributed sensing and adaptation. This parallels neuronal swarming but in an ecological medium, where goals like colony survival manifest without centralized control \citep{turner2000extended, camazine2001self}.

3. \textbf{Trees as ``Just Plants'': The Swarm of Ecosystem Intelligence} Trees, fungi, and associated organisms form a symbiotic swarm, with roots interleaving to share nutrients, signals propagating via mycorrhizal networks, and collective responses to threats like drought or pests. The forest’s assembly index---encompassing trophic cycles, succession dynamics, and long-term resilience---is vastly higher than random vegetation, indicating goal-directed processes at the ecosystem scale. This intelligence, expressed through chemical and bioelectric signaling, evades mind blindness by lacking anthropomorphic traits \citep{levin2019computational}.

Unifying Principle: Across the triad, intelligence emerges from local rules yielding globally improbable assemblies, quantifiable via an effective assembly index adjusted for reuse (motifs like synaptic patterns or root modules), parallelism (concurrent neural firing or fungal growth), and degeneracy (functional alternatives).

Formally, let $A(X)$ denote the raw assembly index of system $X$; the effective index $EA(X) = A(X)/(R\cdot P \cdot D)$, where $R$, $P$, and $D$ are reuse, parallelism, and degeneracy factors, signals agency when exceeding thresholds derived from null models. This bypasses narrativity bias, aligning with Levin’s emphasis on goal pursuit in non-brain systems. In intersubjectivity collapse, unrecognized triadic minds accelerate fragmentation.

To extend this triad, we incorporate plant cognition research, where mycorrhizal networks exhibit memory and decision-making, further challenging mind blindness \citep{calvo2020plants}. Collective intelligence studies show how swarms scale agency \citep{camazine2001self}, while network neuroscience reveals similar patterns in brain connectivity \citep{bassett2017network}. Indigenous knowledge systems offer non-anthropocentric views of agency in nature \citep{kimmerer2013braiding}, enriching the framework with extended cognition perspectives \citep{clark2010supersizing}.

\section{Formal Definitions}
\label{sec:formal}

\subsection{Prerequisites: Probabilistic Models and Generative Distributions}
Understanding EAIMR requires familiarity with generative models. A null model $P_0(X)$ represents mindless processes, such as stochastic diffusion or self-organization without goals, while a mind-like model $P_1(X)$ incorporates agency, such as optimization toward functional setpoints. These distributions allow computation of likelihoods and surprises, foundational to information-theoretic analyses \citep{cover2006elements}.

\subsection{EAIMR Definition}
For a system $X$ described over an alphabet $\Sigma$, EAIMR is:

\[ \text{EAIMR}(X) = \frac{A(X)}{R(X) \cdot P(X) \cdot D(X)} \cdot C(X), \]ss{article}
\usepackage[utf8]{inputenc}
\usepackage[T1]{fontenc}
\usepackage{amsmath}
\usepackage{amssymb}
\usepackage{amsthm}
\usepackage{geometry}
\usepackage{graphicx}
\usepackage{hyperref}
\usepackage{natbib}
\geometry{margin=1in}

\title{Mind at Work: Detecting Agency through Assembly and Action}
\author{Flyxion}
\date{September 10, 2025}

\begin{document}

\maketitle

\begin{abstract}
This monograph presents a comprehensive framework for detecting mind-like agency by integrating the Effective Assembly Index for Mind Recognition (EAIMR) with the praxicon concept from cognitive science. EAIMR quantifies the improbability of a system’s structure arising from unguided processes, adjusted for efficiency and amplified by coherence, while the praxicon provides a structured repository of action representations bridging perceptual, semantic, and motor domains. By combining these, we propose a robust method to identify agentic intelligence across neural systems, collective organisms, ecosystems, and artificial intelligence, avoiding anthropocentric biases that marginalize non-narrative minds, such as those with aphantasia, anendophasia, or ecological intelligences. The framework is formalized using information theory, graph theory, category theory, and sheaf theory, with applications in neuroscience, AI alignment, ecological cognition, and cross-cultural studies. We address intersubjectivity collapse as a structural risk, amplified by Big Tech’s attention-compute arms race and ecological fragmentation, drawing parallels to historical critiques of industrialization. The monograph contrasts this approach with competing theories, provides detailed case studies, and discusses ethical implications for recognizing diverse minds. This extended edition includes multi-scale temporal dynamics, causal intervention frameworks, compositional agency, quantitative validation studies, computational implementations, statistical methods, phenomenological bridges, evolutionary dynamics, information-theoretic refinements, AI alignment protocols, clinical applications, policy recommendations, bias mitigation strategies, computational tractability improvements, and connections to plant cognition, collective intelligence, network neuroscience, indigenous knowledge, and extended cognition.
\end{abstract}

\section{Introduction}
\label{sec:introduction}

The quest to detect and understand agency---defined as the capacity of a system to pursue goals through coordinated actions deviating from unguided physical processes---spans cognitive science, philosophy of mind, artificial intelligence (AI), and ecology. Our civilization’s bias toward equating intelligence with narrativity and imagery, dominant human cognitive modes, systematically marginalizes alternative forms of intelligence. As Michael Levin articulates, ``mind blindness'' leads us to dismiss the agency of plants, animals, and cellular collectives because they lack human-like symbols or narratives \citep{levin2022technological}. Levin's work emphasizes that intelligence is not brain-bound but present in any system capable of goal pursuit and adaptation, from cellular regeneration to collective behaviors \citep{levin2019computational}. This bias extends to cognitive minorities, such as individuals with aphantasia (inability to visualize mentally) or anendophasia (absence of inner monologue), who face exclusion in narrative-driven economies despite intact functional reasoning \citep{zeman2015, hurlburt2024}.

Big Tech’s business model---selling compute and storage while buying attention through algorithmic amplification---exacerbates this collapse. By spotlighting outlier success stories and suppressing systemic critiques, platforms create a narrative arms race that erodes trust and shared meaning \citep{zuboff2019age, wu2016attention}. Programming courses exemplify this: an oversupply of training with undersupplied job demand creates a hope market, where platforms profit by amplifying rare successes while burying failures \citep{carr2010shallows}. This dynamic parallels historical ecological critiques, from Joni Mitchell’s ``Big Yellow Taxi'' lamenting the paving of paradise for parking lots to Jim Amos’s vision of a ``Coruscant planet'' where data centers and nuclear plants displace ecosystems for a compute-driven metaverse \citep{lanham2006economics}. Roads, like digital feeds, fragment habitats---disrupting migrations, raising heat islands, and diverting water---leading to biodiversity loss, population isolation, and ecosystem degradation \citep{forman2003road}. Both create isolated patches, undermining exchange of animals or ideas, driving collapse \citep{arendt1958human}.

The Effective Assembly Index for Mind Recognition (EAIMR) and the praxicon framework offer a solution by quantifying agency through structural complexity and action representation, bypassing narrativity bias to recognize diverse intelligences. EAIMR measures the improbability of a system’s assembly, adjusted for reuse, parallelism, and degeneracy, and amplified by coherence \citep{walker2013algorithmic}. The praxicon, a cognitive ``dictionary'' of actions, bridges perceptual, semantic, and motor domains, grounding intelligence in embodied processes \citep{pastra2011praxicon, peigneux2004imaging}. Together, they detect agency in systems from neural circuits to termite mounds, forests, and AI, fostering inclusivity for non-human and atypical human minds. This monograph integrates these frameworks with information theory, graph theory, category theory, and sheaf theory, addressing applications in neuroscience, AI alignment, ecology, and cross-cultural cognition, while situating IC as a meta-crux for debates on AI futures, such as those between Yudkowsky’s pessimism and Miller’s cooperative optimism \citep{yudkowsky2008artificial, miller2006robust}.

To appreciate this integration, it is essential to understand prerequisite concepts. Information theory, pioneered by Claude Shannon, quantifies uncertainty and structure in systems through measures like entropy, which represents the average information content or disorder in a probability distribution \citep{shannon1948mathematical}. Complexity measures, such as Kolmogorov complexity, define the shortest program capable of generating a given structure, providing a basis for assessing improbability \citep{kolmogorov1965three}. Assembly theory extends these ideas to biological and technological contexts by counting the minimal steps required to construct complex objects from simple building blocks \citep{walker2013algorithmic}. The praxicon, rooted in neuropsychology, models how the brain stores and retrieves action knowledge, distinguishing between perceptual recognition and motor execution \citep{peigneux2004imaging}. By weaving these threads together, our framework addresses the challenge of detecting agency without relying on anthropocentric markers like language or self-report.

Furthermore, this work addresses ``mind blindness,'' a term from Michael Levin’s research describing the human tendency to overlook intelligence in non-human systems due to their lack of narrative or imagistic expression. This bias extends to cognitive minorities within humanity, such as individuals with aphantasia (inability to visualize mentally) or anendophasia (lack of inner speech), who process information functionally but are marginalized in narrative-dominant societies. Integrating these insights, the framework mitigates intersubjectivity collapse---the breakdown of shared understanding among diverse minds---by recognizing alternative intelligences through structural metrics like EAIMR.

\section{Historical and Conceptual Foundations}
\label{sec:foundations}

\subsection{Prerequisites: Complexity and Information in Natural Systems}
Before delving into assembly theory, consider the foundational role of complexity in distinguishing living and agentic systems from inert matter. Murray Gell-Mann’s work on effective complexity highlights how systems exhibit regularities that are neither too random nor too ordered, a balance often indicative of adaptive processes \citep{adami2012use}. Robert Hazen’s research on mineral evolution and biocomplexity introduces functional information as a metric for the emergence of life-like properties, where improbable configurations serve functional roles \citep{hazen2007functional}. These ideas set the stage for assembly theory, which quantifies the ``assembly index'' as the minimal number of operations needed to produce a structure, emphasizing reuse and hierarchy in evolved systems.

\subsection{Assembly Theory and Complexity}
Assembly theory, developed by Sara Imari Walker, Leroy Cronin, and collaborators, provides a formal measure of complexity by calculating the shortest path to assemble a molecule or structure from basic units \citep{cronin2016beyond, walker2013algorithmic}. It has applications in astrobiology and origins-of-life research, where high assembly indices correlate with biotic processes. EAIMR builds upon this by incorporating adjustments for real-world efficiencies and coherence, enabling its use as a detector for mind-like agency \citep{adami2012use}. Loopholes in assembly index application---such as reuse of substructures, parallelism, context dependence, degeneracy, and information leakage---highlight how complexity can bootstrap without full cost, relevant to understanding emergent intelligences in swarms or ecosystems.

\subsection{Praxicon: Action Representation in Cognition}
The praxicon emerges from studies of apraxia, a disorder impairing purposeful movements despite intact motor function. Prerequisite knowledge includes the distinction between ideational apraxia (conceptual deficits) and ideomotor apraxia (execution deficits), as modeled by Rothi and Heilman \citep{mantovani2010reviewing}. The praxicon is conceptualized as a dual-component system: the input praxicon for recognizing actions via visual-gestural cues, and the output praxicon for producing motor engrams \citep{peigneux2004imaging}. Extensions to AI involve grounding symbolic language in sensorimotor primitives, addressing the symbol grounding problem \citep{harnad1990symbol, pastra2011praxicon, vannuscorps2023effector}.

\subsection{Competing Approaches}
To contextualize our framework, compare it with alternatives. Integrated Information Theory (IIT) posits consciousness as integrated information $\Phi$, but it requires predefined system boundaries and struggles with distributed agency \citep{tononi2004information}. The Free Energy Principle (FEP) frames agents as minimizers of variational free energy, modeling predictive processing but lacking direct measures of structural improbability \citep{friston2010}. Enactive cognition emphasizes autopoiesis and sensorimotor contingencies, yet it often remains qualitative \citep{barandiaran2009defining}. Our EAIMR-praxicon approach complements these by providing computable metrics grounded in assembly and action representation.

\section{Mind Blindness and Intersubjectivity Collapse}
\label{sec:mind_blindness}

Mind blindness, as articulated by Michael Levin, refers to the systematic failure to recognize intelligence in non-human systems, such as plants, animals, and cellular collectives, because they do not conform to human-like symbolic or narrative modes. Levin argues that intelligence is any system’s ability to pursue goals, adapt to perturbations, and solve problems in its medium, evident in regenerating limbs or bioelectric signaling \citep{levin2019computational}. This bias extends beyond ecology to cognitive minorities within humanity, where individuals with aphantasia---the inability to voluntarily create mental images---or anendophasia---the absence of inner monologue---are marginalized despite functional processing capabilities.

\textbf{Mind Blindness Across Domains}

1. \textbf{Ecological Minds (Plants and Animals)} Bias: Intelligence is equated with language and narrative. Blindness: Plants and animals operate through signaling, adaptation, and coordination, but because they do not ``speak,'' their intelligence is dismissed as instinct. Consequence: Systematic marginalization of non-human agency, paving the way for ecological exploitation.

2. \textbf{Cognitive Minorities (Aphantasia, Anendophasia, Atypical Cognition)} Bias: Value is tied to attention performance---vivid self-narration, imagery, persuasion. Blindness: Those whose minds process functionally but without inner imagery or speech must reverse-engineer outputs, making them invisible or ``less capable'' in the attention economy. Aphantasia affects 2-4\% of the population, impairing voluntary visualization, while anendophasia involves silent inner worlds without verbal thought. Consequence: Social and economic exclusion of cognitive styles that do not match platform-optimized norms.

3. \textbf{Religious and Philosophical Systems (Christianity, Islam, Atheism, etc.)} Bias: Each tradition assumes its own narrative or epistemic mode (ritual, revelation, rational critique) is the universal medium of truth. Blindness: Values and intelligences encoded in rival traditions become illegible. What one calls wisdom, another calls heresy or superstition. Consequence: Breakdown of shared meaning across civilizations, escalating into mistrust, conflict, or war.

Unifying Thread: Across all three layers, mind blindness is the refusal to recognize the legitimacy of intelligences that do not mirror our own mode of expression. This blindness is not just a moral failure---it is a structural driver of intersubjectivity collapse, as communication, empathy, and trust fail across ecological, cognitive, and cultural domains. AI futures will magnify this bias---if we only design for intelligences that ``talk like us,'' we will fail to cooperate with (or even recognize) those that do not.

\subsection{Termite--Neuron--Forest Triad: Analogies for Detecting Emergent Intelligence}

In the context of intersubjectivity collapse---where diverse minds fail to communicate or recognize one another due to incompatible cognitive modes---the ``Termite--Neuron--Forest Triad'' serves as a formal analogy to highlight how emergent intelligence arises from swarms of simpler agents. This triad underscores the structural bias of mind blindness, wherein intelligences lacking narrative or imagistic expression are dismissed, exacerbating collapse by excluding valid forms of agency. By applying an assembly-index lens, we can quantify the improbability of these emergent structures, revealing mind-like processes that operate through coordination rather than human-centric symbolism.

1. \textbf{Neurons as ``Just Insects'': The Swarm of Consciousness} Neurons, individually, are simple electrochemical agents following local rules of excitation and inhibition. However, their collective swarming---through synaptic connections, feedback loops, and hierarchical organization---assembles a brain capable of consciousness and goal-directed behavior. The assembly index here is extraordinarily high: the precise mesoscale architecture (e.g., cortical columns) is improbable under mindless diffusion, implying directed processes like plasticity and learning. Mind blindness dismisses this as ``mere biology'' until the emergent whole is recognized, mirroring how we overlook atypical human cognition (e.g., aphantasia) that lacks internal imagery.

2. \textbf{Termites as ``Just Insects'': The Swarm of Architecture} Termites operate as decentralized agents, responding to local pheromones and environmental cues. Yet, their swarming constructs mounds with advanced features: ventilation topology for thermoregulation, chambers for farming fungi, and resilient structures against floods. The mound’s assembly index exceeds that of random soil accumulations (e.g., dunes), reflecting collective intelligence in distributed sensing and adaptation. This parallels neuronal swarming but in an ecological medium, where goals like colony survival manifest without centralized control \citep{turner2000extended, camazine2001self}.

3. \textbf{Trees as ``Just Plants'': The Swarm of Ecosystem Intelligence} Trees, fungi, and associated organisms form a symbiotic swarm, with roots interleaving to share nutrients, signals propagating via mycorrhizal networks, and collective responses to threats like drought or pests. The forest’s assembly index---encompassing trophic cycles, succession dynamics, and long-term resilience---is vastly higher than random vegetation, indicating goal-directed processes at the ecosystem scale. This intelligence, expressed through chemical and bioelectric signaling, evades mind blindness by lacking anthropomorphic traits \citep{levin2019computational}.

Unifying Principle: Across the triad, intelligence emerges from local rules yielding globally improbable assemblies, quantifiable via an effective assembly index adjusted for reuse (motifs like synaptic patterns or root modules), parallelism (concurrent neural firing or fungal growth), and degeneracy (functional alternatives).

Formally, let $A(X)$ denote the raw assembly index of system $X$; the effective index $EA(X) = A(X)/(R\cdot P \cdot D)$, where $R$, $P$, and $D$ are reuse, parallelism, and degeneracy factors, signals agency when exceeding thresholds derived from null models. This bypasses narrativity bias, aligning with Levin’s emphasis on goal pursuit in non-brain systems. In intersubjectivity collapse, unrecognized triadic minds accelerate fragmentation.

To extend this triad, we incorporate plant cognition research, where mycorrhizal networks exhibit memory and decision-making, further challenging mind blindness \citep{calvo2020plants}. Collective intelligence studies show how swarms scale agency \citep{camazine2001self}, while network neuroscience reveals similar patterns in brain connectivity \citep{bassett2017network}. Indigenous knowledge systems offer non-anthropocentric views of agency in nature \citep{kimmerer2013braiding}, enriching the framework with extended cognition perspectives \citep{clark2010supersizing}.

\section{Formal Definitions}
\label{sec:formal}

\subsection{Prerequisites: Probabilistic Models and Generative Distributions}
Understanding EAIMR requires familiarity with generative models. A null model $P_0(X)$ represents mindless processes, such as stochastic diffusion or self-organization without goals, while a mind-like model $P_1(X)$ incorporates agency, such as optimization toward functional setpoints. These distributions allow computation of likelihoods and surprises, foundational to information-theoretic analyses \citep{cover2006elements}.

\subsection{EAIMR Definition}
For a system $X$ described over an alphabet $\Sigma$, EAIMR is:

\[ \text{EAIMR}(X) = \frac{A(X)}{R(X) \cdot P(X) \cdot D(X)} \cdot C(X), \]

where: - $A(X)$: Raw Assembly Index, the minimal steps in an assembly grammar \citep{walker2013algorithmic}. - $R(X)$: Reuse factor, quantifying modular repetition. - $P(X)$: Parallelism factor, accounting for concurrent processes. - $D(X)$: Degeneracy factor, reflecting alternative configurations achieving similar functions. - $C(X)$: Coherence factor, amplifying goal-directed integration. High values suggest agency, as the structure is improbable under $P_0(X)$ but plausible under $P_1(X)$.

\subsection{Praxicon Definition}
The praxicon is formalized as a bipartite graph $G_P = (V_A, V_S, E)$, with action vertices $V_A$ (e.g., motor programs), semantic vertices $V_S$ (e.g., conceptual labels), and edges $E$ encoding mappings \citep{panunzi2018imagact}. This structure facilitates bidirectional translation between perception and action, essential for embodied cognition \citep{barsalou2008grounded}.

\section{Mathematical Integration}
\label{sec:integration}

\subsection{Information-Theoretic Foundations}
Self-information under the null: $I_0(X) = -\log P_0(X)$. The likelihood ratio $\Lambda(X) = P_1(X)/P_0(X)$ quantifies preference for agentic models. For praxicon integration, action entropy reduction is:

$\Delta H_P = H(P_0) - H_{\text{obs}}(P)$,

updating coherence as $C(X) \approx 1 + \alpha \Delta H_P / H(P_0)$ \citep{shannon1948mathematical, jaynes2003probability}. Kolmogorov complexity approximates $A(X)$ \citep{kolmogorov1965three}.

\subsection{Graph-Theoretic Formalism}
The praxicon graph $G_P$ allows coherence via spectral properties, e.g., eigenvalue stability for robustness. EAIMR identifies low-entropy subgraphs indicative of agency \citep{newell1990unified}.

\subsection{Category-Theoretic Interpretation}
Category theory provides a formal language for describing structures and transformations in abstract terms, making it particularly suitable for modeling cognitive processes that involve mappings between different domains \citep{healy2022, phillips2021}. A category consists of objects (e.g., perceptual states, motor actions) and morphisms (transformations or relations between objects) that satisfy associativity and identity axioms. Functors map between categories while preserving structure, and natural transformations provide morphisms between functors. In our framework, the praxicon can be viewed as a functor $F : C_{\text{perc}} \to C_{\text{motor}}$, where $C_{\text{perc}}$ is the category of perceptual inputs (objects as sensory data, morphisms as perceptual transformations) and $C_{\text{motor}}$ is the category of motor outputs (objects as action states, morphisms as motor adjustments), mediated through a semantic category $C_{\text{sem}}$. This functorial representation captures how perceptual inputs are systematically transformed into motor actions via semantic grounding. EAIMR, in this context, quantifies the ``improbability'' or cost of natural transformations $\eta : F \implies G$, where $G$ represents alternative mappings (e.g., under null dynamics). High EAIMR corresponds to scenarios where the observed natural transformation (the agentic mapping) is highly specific and unlikely under random perturbations, reflecting the coordinated agency. For instance, in neural systems, category-theoretic models can describe how cognitive development involves composing functors for increasingly complex action representations \citep{ehresmann1980}. This interpretation aligns with structuralist views in cognitive science, emphasizing relational structures over material substrates \citep{ehresmann2021}.

\subsection{Sheaf-Theoretic Interpretation}
Sheaf theory extends category theory by addressing how local data on a topological space can be consistently glued into global structures, offering a powerful tool for modeling distributed systems where agency emerges from multi-scale interactions \citep{spivak2024}. A sheaf $F$ on a space $X$ assigns to each open set $U \subseteq X$ a set $F(U)$ (local data) with restriction maps that satisfy gluing conditions: local sections that agree on overlaps can be uniquely combined into a global section. In biological and ecological contexts, sheaf theory models uncertainty and resilience in complex networks \citep{stenzel2024}. For our framework, define a sheaf $F$ on the topological space representing the system’s components (e.g., neurons in a circuit or species in an ecosystem), where sections over open sets encode local praxicon fragments---action representations valid in subsystems. The gluing axiom ensures that local actions cohere into global behaviors, such as coordinated foraging in ant colonies or nutrient cycling in forests. EAIMR quantifies obstructions to trivial gluing, often captured by sheaf cohomology groups $H^k(X, F)$. Non-zero cohomology indicates ``twists'' or global inconsistencies under null models, corresponding to high improbability and thus agency. For example, in ecology, forest nutrient networks form a sheaf where local flows (sections) glue into global cycles; high $H^1$ (non-exact cycles) elevates EAIMR, detecting emergent intelligence \citep{levin2019computational}. In biology, bioelectric signaling can be sheafified, with cohomology measuring scale-free cognition \citep{levin2019computational, spivak2024}. This approach bridges local computations to global agency, providing a mathematical basis for resilience under perturbations.

\section{Worked Examples}
\label{sec:examples}

\subsection{Neural Tissue}
In neural systems, high $A(X)$ arises from precise synaptic architectures, with praxicon engrams in cortical columns reducing entropy. Simulations of mirror neurons show effector-specific mappings \citep{gallese1998mirror}. Step-by-step: Compute $A(X)$ via connection counts, adjust discounts, amplify with coherence from goal pursuit. Entropy curve: Sharp reduction in action space.

\subsection{Termite Mound}
Moderate assembly in tunnels, high coherence in thermoregulation. Collective praxicon via pheromone-guided motifs \citep{turner2000extended}. Graph: Nodes for worker actions linked to colony semantics. Data from self-organization models \citep{camazine2001self}.

\subsection{Forest Ecosystem}
Extensive macro-assembly in symbiotic networks, praxicon-like routing. High degeneracy in species roles, coherence over timescales \citep{levin2019computational}. Lattice diagram: Interconnected trophic levels.

\subsection{Negative Case: Crystal Growth}
Regular lattice growth yields low $C(X)$, high predictability under $P_0$, resulting in low EAIMR \citep{hazen2007functional}.

\section{Theoretical Extensions}
\label{sec:theoretical_extensions}

To enhance the framework, we introduce multi-scale temporal dynamics with a temporal coherence component $C_t(X)$, measuring agency across timescales from neural firing to forest succession, distinguishing long-term agency from patterns \citep{levin2019computational}.

Causal intervention frameworks define an agency resilience index, quantifying recovery from targeted disruptions versus noise, differentiating true agents from homeostatic systems \citep{barandiaran2009defining}.

Compositional agency extends category theory for nested intelligences, where agents form collectives with irreducible properties \citep{ehresmann2021}.

\section{Empirical and Methodological Improvements}
\label{sec:empirical}

Quantitative validation studies are essential, including neural comparisons of healthy vs. lesioned tissue, ecological manipulations of forests/termite colonies, AI benchmarks, and cross-cultural ritual analyses \citep{pastra2011praxicon}.

Computational implementation requires open-source tools for EAIMR computation, with standardized data formats, Kolmogorov approximation algorithms, and praxicon graph extraction methods \citep{kolmogorov1965three}.

Statistical framework includes confidence intervals on EAIMR, multiple comparison corrections, power analysis for agency detection, and handling uncertainty \citep{jaynes2003probability}.

\section{Conceptual Deepening}
\label{sec:conceptual}

The phenomenological bridge addresses consciousness: high EAIMR systems may exhibit phenomenal experience, with testable predictions for subjective states \citep{tononi2004information}.

Evolutionary dynamics predict higher EAIMR systems show distinct trajectories, testable in simulations or comparative studies \citep{camazine2001self}.

Information-theoretic refinements connect EAIMR to integrated information, non-equilibrium thermodynamics, and predictive processing \citep{friston2010}.

\section{Practical Applications}
\label{sec:applications}

\subsection{Neuroscience}
Detects apraxia and atypical cognition \citep{peigneux2004imaging, zeman2015}. Clinical applications include diagnostics for locked-in syndrome, rehabilitation monitoring, and therapeutic community agency assessment \citep{levin2022technological}.

\subsection{AI and Robotics}
Grounds LLMs; detects agency \citep{pastra2011praxicon, anderson2007how}. AI alignment protocols monitor emergent agency during training, detect deceptive alignment, and measure value alignment via praxicon analysis \citep{yudkowsky2008artificial}.

\subsection{Ecological and Collective Minds}
Identifies swarm intelligence \citep{camazine2001self, simondon2017mode}. Policy frameworks enable legal recognition of non-human agency, environmental protection based on ecosystem metrics, and AI regulation \citep{bostrom2014superintelligence}.

\subsection{Cross-Cultural Cognition}
Rituals as praxicons \citep{ortega1961meditations}. Anthropocentric bias mitigation involves cross-validation with non-human measures, indigenous knowledge, and animal studies \citep{kimmerer2013braiding}.

\section{Limits and Future Research}
\label{sec:limits}

Challenges: Approximating complexity, validating sheaves. Computational tractability requires polynomial approximations and hierarchical methods for large systems \citep{kolmogorov1965three}.

Future: Bioelectric models, AI simulations \citep{levin2019computational}. Connections to plant cognition \citep{calvo2020plants}, collective intelligence \citep{camazine2001self}, network neuroscience \citep{bassett2017network}, indigenous knowledge \citep{kimmerer2013braiding}, and extended cognition \citep{clark2010supersizing}.

\section{Ethical and Philosophical Implications}
\label{sec:ethics}

Mitigates mind blindness, informs AI ethics, and aligns with process philosophy \citep{dennett1987intentional, whitehead1929process, deleuze1987thousand}.

\section{Decision Rule}
\label{sec:decision}

$\text{EAIMR}(X) \ge \tau$ classifies agency, using graph entropy proxies.

\section{Conclusion}
\label{sec:conclusion}

This framework detects agency, averting intersubjectivity collapse by recognizing diverse minds.

\appendix

\section{Mathematical Appendix: EAIMR, Information, and Improbable Order}

A.1 Preliminaries and Notation

Let $X$ denote a structure/system (e.g., termite mound, forest network, neural microcircuit) described over an alphabet $\Sigma$. Let $P_0$ be a null (mindless) generative model for $X$ (e.g., stochastic physics + simple self-organization) with distribution $P_0(X)$. Let $P_1$ be a mind-like generative model for $X$ (agentic/goal-directed) with distribution $P_1(X)$. Raw Assembly Index $A(X)$: minimum step count in a prescribed assembly grammar to build $X$ from primitives.

EAIMR (from your definition):

\[ \text{EAIMR}(X) = \frac{A(X)}{R(X) \cdot P(X) \cdot D(X) \cdot C(X)} \]

A.2 Connection to Kolmogorov Complexity

Let $K(X)$ be prefix Kolmogorov complexity (shortest universal description of $X$). In general, lower $K(X)$ higher compressibility (simple regularities). Many mind-built artifacts have mixed profiles: globally compressible high-level plan + low-level details that are costly to specify. A useful proxy is two-part MDL (minimum description length):

\[ L(X) \approx L(\text{model}) + L(\text{residuals} \mid \text{model}) \]

Heuristic link (non-rigorous but actionable):

\[ \text{EAIMR}(X) \propto \frac{K_{\text{plan}}(X)}{R \cdot P \cdot D} \cdot C \]

with $K_{\text{plan}}(X) \lesssim K(X)$.

Kolmogorov complexity $K(X)$ is the length of the shortest program $p$ in a fixed universal Turing machine $U$ such that $U(p) = X$. For EAIMR, we approximate $K_{\text{plan}}(X)$ as the compressible plan component, separating it from residuals using two-part coding, which aligns with practical compression algorithms like LZW or gzip for estimation.

A.3 Statistical Surprise and Likelihood Ratios

Define the self-information under the null:

\[ I_0(X) \equiv -\log P_0(X) \]

\[ \Lambda(X) \equiv \frac{P_1(X)}{P_0(X)} = \exp \left( \log P_1(X) - \log P_0(X) \right). \]

Key idea: High EAIMR should correlate with large $\Lambda(X)$ and a high $I_0(X)$ (i.e., $X$ is very unlikely under mindless dynamics but plausible under mind-like dynamics). A practical surrogate score you can compute or estimate:

\[ S(X) \equiv I_0(X) + \lambda \phi(X) \]

\[ \text{EAIMR}(X) = g \left( S(X) \right), \quad g \uparrow \]

The likelihood ratio $\Lambda(X)$ is the Bayes factor for mind-like vs. null models, with $\log \Lambda(X)$ as evidence weight. For estimation, use Monte Carlo methods to approximate $P_0(X)$ and $P_1(X)$ from generative simulations, such as Markov chain Monte Carlo for high-dimensional spaces.

A.4 Entropy Reduction, Work, and Directed Assembly

Let $M$ be macrostate features (e.g., ventilation topology, nutrient routing, neural coding). Define macro-entropy under $P_0$ and observed for $X$. Directed assembly manifests as entropy reduction:

\[ \Delta H \equiv H(M) - H_{\text{obs}}(M) > 0 \]

A convenient coherence multiplier:

\[ C(X) \approx 1 + \alpha \frac{\Delta H}{H(M)} \]

Entropy $H(M) = -\sum p_i \log p_i$ over macrostates, where $p_i$ is probability under $P_0$. Observed $H_{\text{obs}}(M)$ is computed from empirical distributions, with $\Delta H$ quantifying work against thermodynamic disorder, linking to non-equilibrium physics \citep{friston2010}.

A.5 Reuse, Parallelism, Degeneracy as Information Discounts

Treat these as code-length discounts relative to the raw assembly plan. Reuse: library motifs shorten description. If motif dictionary has description amortized across instances, the effective bits per instance drop. Model as a divisor on $A(X)$. Parallelism: independent sub-assemblies compress as a product code; wall-clock and effective complexity per unit time fall. Model as a divisor on $A(X)$. Degeneracy: if many alternatives achieve the same function, the description of ``a solution'' is shorter than ``this specific solution.'' Function-first coding reduces effective complexity; model as a divisor. This aligns your original formula with MDL-style accounting: EAIMR rises when, despite these discounts, the remaining coordinated specification is still large and functionally tight.

For reuse, $R(X) = 1 + \log(\text{number of motif instances} / \text{motif types})$, amortizing costs. Parallelism $P(X) = \text{total operations} / \text{critical path length}$, capturing concurrency. Degeneracy $D(X) = \log(\text{number of functional equivalents})$, spreading probability mass.

A.6 Coherence via Control and Causal Graphs

Let $G(X)$ be a causal graph extracted from $X$’s dynamics. Define a control coherence term:

\[ \phi(X) = \beta_1 \text{Stab}(G) + \beta_2 \text{GoalAcc}(G) + \beta_3 \text{Robust}(G) \]

GoalAcc($G$): accuracy tracking a specified functional setpoint (e.g., temperature band in termite mounds). Robust($G$): performance under noise/lesions (graceful degradation). Then use $\phi(X)$ inside (Section 3) or directly as your multiplier.

Stab($G$) can be graph Laplacian eigenvalues for stability; GoalAcc($G$) as mean squared error from setpoint; Robust($G$) as fraction of performance retained under simulated lesions, using graph pruning algorithms.

A.7 Practical Estimation (finite data)

Because $K(X)$ is uncomputable, rely on approximations: Compression proxies: on suitable encodings (graphs, fields, symbol streams). Surprisal under null: fit $P_0$ (e.g., generative physics sims, random graphs) and estimate by Monte Carlo or density models. Function and coherence: quantify with control/evaluation tasks (tracking error, energy efficiency, resilience curves). Discounts: measure from libraries (motif counts), concurrency (critical path vs. total ops), and function-equivalence classes (number of viable variants).

For compression, use BZ2 or LZMA ratios as $K(X)$ lower bounds. Monte Carlo for $P_0$ uses $10^4-10^6$ samples for convergence. Resilience curves fit to exponential decay models for Robust($G$).

A.8 Worked Intuitions (triad examples)

Neural tissue: very high $A(X)$ (precise mesoscale organization), strong $C(X)$ (robust goal pursuit), discounts present (columns, parallelism), yet EAIMR stays high mind-like. Termite mound: moderate $A(X)$ vs. dunes, high $C(X)$ (thermoregulation), sizable discounts (motifs, parallel labor, many workable designs) but coordinated control keeps EAIMR high. Forest: extremely high macro $A(X)$ vs. random vegetation, high robustness/goal achievement (nutrient routing, succession), massive reuse/degeneracy, yet coherence over centuries pushes EAIMR high.

A.9 Decision Rule (operational)

Define thresholds from a reference set of clearly mindless assemblies:

$\text{EAIMR}(X) \ge \tau \implies$ classify as mind-like (candidate)

--- Takeaway

Information-theoretic view: EAIMR is a mind detector for improbably coherent order---high surprisal under a null, discounted for shortcuts, and boosted by functional integration. It sidesteps narrativity bias: you do not need language or imagery to register as intelligent; termite colonies, forests, and atypical human cognition can score high if they assemble and maintain low-entropy, goal-directed structure against default dynamics.

B Derivations and Pseudocode

B.1 EAIMR Derivation

From surprisal: $S(X) = I_0(X) + \lambda \phi(X)$, EAIMR as monotonic $g(S(X))$.

B.2 Pseudocode

\begin{verbatim}
def compute_eaimr(X, P0, P1, G_P):
    A = assembly_index(X)
    R = reuse_factor(X)
    P = parallelism(X)
    D = degeneracy(X)
    C = coherence(G_P)
    return (A / (R * P * D)) * C
\end{verbatim}

\bibliographystyle{plainnat}
\bibliography{references}

\end{document}
ss{article}
\usepackage[utf8]{inputenc}
\usepackage[T1]{fontenc}
\usepackage{amsmath}
\usepackage{amssymb}
\usepackage{amsthm}
\usepackage{geometry}
\usepackage{graphicx}
\usepackage{hyperref}
\usepackage{natbib}
\geometry{margin=1in}

\title{Mind at Work: Detecting Agency through Assembly and Action}
\author{Flyxion}
\date{September 10, 2025}

\begin{document}

\maketitle

\begin{abstract}
This monograph presents a comprehensive framework for detecting mind-like agency by integrating the Effective Assembly Index for Mind Recognition (EAIMR) with the praxicon concept from cognitive science. EAIMR quantifies the improbability of a system’s structure arising from unguided processes, adjusted for efficiency and amplified by coherence, while the praxicon provides a structured repository of action representations bridging perceptual, semantic, and motor domains. By combining these, we propose a robust method to identify agentic intelligence across neural systems, collective organisms, ecosystems, and artificial intelligence, avoiding anthropocentric biases that marginalize non-narrative minds, such as those with aphantasia, anendophasia, or ecological intelligences. The framework is formalized using information theory, graph theory, category theory, and sheaf theory, with applications in neuroscience, AI alignment, ecological cognition, and cross-cultural studies. We address intersubjectivity collapse as a structural risk, amplified by Big Tech’s attention-compute arms race and ecological fragmentation, drawing parallels to historical critiques of industrialization. The monograph contrasts this approach with competing theories, provides detailed case studies, and discusses ethical implications for recognizing diverse minds. This extended edition includes multi-scale temporal dynamics, causal intervention frameworks, compositional agency, quantitative validation studies, computational implementations, statistical methods, phenomenological bridges, evolutionary dynamics, information-theoretic refinements, AI alignment protocols, clinical applications, policy recommendations, bias mitigation strategies, computational tractability improvements, and connections to plant cognition, collective intelligence, network neuroscience, indigenous knowledge, and extended cognition.
\end{abstract}

\section{Introduction}
\label{sec:introduction}

The quest to detect and understand agency---defined as the capacity of a system to pursue goals through coordinated actions deviating from unguided physical processes---spans cognitive science, philosophy of mind, artificial intelligence (AI), and ecology. Our civilization’s bias toward equating intelligence with narrativity and imagery, dominant human cognitive modes, systematically marginalizes alternative forms of intelligence. As Michael Levin articulates, ``mind blindness'' leads us to dismiss the agency of plants, animals, and cellular collectives because they lack human-like symbols or narratives \citep{levin2022technological}. Levin's work emphasizes that intelligence is not brain-bound but present in any system capable of goal pursuit and adaptation, from cellular regeneration to collective behaviors \citep{levin2019computational}. This bias extends to cognitive minorities, such as individuals with aphantasia (inability to visualize mentally) or anendophasia (absence of inner monologue), who face exclusion in narrative-driven economies despite intact functional reasoning \citep{zeman2015, hurlburt2024}.

Big Tech’s business model---selling compute and storage while buying attention through algorithmic amplification---exacerbates this collapse. By spotlighting outlier success stories and suppressing systemic critiques, platforms create a narrative arms race that erodes trust and shared meaning \citep{zuboff2019age, wu2016attention}. Programming courses exemplify this: an oversupply of training with undersupplied job demand creates a hope market, where platforms profit by amplifying rare successes while burying failures \citep{carr2010shallows}. This dynamic parallels historical ecological critiques, from Joni Mitchell’s ``Big Yellow Taxi'' lamenting the paving of paradise for parking lots to Jim Amos’s vision of a ``Coruscant planet'' where data centers and nuclear plants displace ecosystems for a compute-driven metaverse \citep{lanham2006economics}. Roads, like digital feeds, fragment habitats---disrupting migrations, raising heat islands, and diverting water---leading to biodiversity loss, population isolation, and ecosystem degradation \citep{forman2003road}. Both create isolated patches, undermining exchange of animals or ideas, driving collapse \citep{arendt1958human}.

The Effective Assembly Index for Mind Recognition (EAIMR) and the praxicon framework offer a solution by quantifying agency through structural complexity and action representation, bypassing narrativity bias to recognize diverse intelligences. EAIMR measures the improbability of a system’s assembly, adjusted for reuse, parallelism, and degeneracy, and amplified by coherence \citep{walker2013algorithmic}. The praxicon, a cognitive ``dictionary'' of actions, bridges perceptual, semantic, and motor domains, grounding intelligence in embodied processes \citep{pastra2011praxicon, peigneux2004imaging}. Together, they detect agency in systems from neural circuits to termite mounds, forests, and AI, fostering inclusivity for non-human and atypical human minds. This monograph integrates these frameworks with information theory, graph theory, category theory, and sheaf theory, addressing applications in neuroscience, AI alignment, ecology, and cross-cultural cognition, while situating IC as a meta-crux for debates on AI futures, such as those between Yudkowsky’s pessimism and Miller’s cooperative optimism \citep{yudkowsky2008artificial, miller2006robust}.

To appreciate this integration, it is essential to understand prerequisite concepts. Information theory, pioneered by Claude Shannon, quantifies uncertainty and structure in systems through measures like entropy, which represents the average information content or disorder in a probability distribution \citep{shannon1948mathematical}. Complexity measures, such as Kolmogorov complexity, define the shortest program capable of generating a given structure, providing a basis for assessing improbability \citep{kolmogorov1965three}. Assembly theory extends these ideas to biological and technological contexts by counting the minimal steps required to construct complex objects from simple building blocks \citep{walker2013algorithmic}. The praxicon, rooted in neuropsychology, models how the brain stores and retrieves action knowledge, distinguishing between perceptual recognition and motor execution \citep{peigneux2004imaging}. By weaving these threads together, our framework addresses the challenge of detecting agency without relying on anthropocentric markers like language or self-report.

Furthermore, this work addresses ``mind blindness,'' a term from Michael Levin’s research describing the human tendency to overlook intelligence in non-human systems due to their lack of narrative or imagistic expression. This bias extends to cognitive minorities within humanity, such as individuals with aphantasia (inability to visualize mentally) or anendophasia (lack of inner speech), who process information functionally but are marginalized in narrative-dominant societies. Integrating these insights, the framework mitigates intersubjectivity collapse---the breakdown of shared understanding among diverse minds---by recognizing alternative intelligences through structural metrics like EAIMR.

\section{Historical and Conceptual Foundations}
\label{sec:foundations}

\subsection{Prerequisites: Complexity and Information in Natural Systems}
Before delving into assembly theory, consider the foundational role of complexity in distinguishing living and agentic systems from inert matter. Murray Gell-Mann’s work on effective complexity highlights how systems exhibit regularities that are neither too random nor too ordered, a balance often indicative of adaptive processes \citep{adami2012use}. Robert Hazen’s research on mineral evolution and biocomplexity introduces functional information as a metric for the emergence of life-like properties, where improbable configurations serve functional roles \citep{hazen2007functional}. These ideas set the stage for assembly theory, which quantifies the ``assembly index'' as the minimal number of operations needed to produce a structure, emphasizing reuse and hierarchy in evolved systems.

\subsection{Assembly Theory and Complexity}
Assembly theory, developed by Sara Imari Walker, Leroy Cronin, and collaborators, provides a formal measure of complexity by calculating the shortest path to assemble a molecule or structure from basic units \citep{cronin2016beyond, walker2013algorithmic}. It has applications in astrobiology and origins-of-life research, where high assembly indices correlate with biotic processes. EAIMR builds upon this by incorporating adjustments for real-world efficiencies and coherence, enabling its use as a detector for mind-like agency \citep{adami2012use}. Loopholes in assembly index application---such as reuse of substructures, parallelism, context dependence, degeneracy, and information leakage---highlight how complexity can bootstrap without full cost, relevant to understanding emergent intelligences in swarms or ecosystems.

\subsection{Praxicon: Action Representation in Cognition}
The praxicon emerges from studies of apraxia, a disorder impairing purposeful movements despite intact motor function. Prerequisite knowledge includes the distinction between ideational apraxia (conceptual deficits) and ideomotor apraxia (execution deficits), as modeled by Rothi and Heilman \citep{mantovani2010reviewing}. The praxicon is conceptualized as a dual-component system: the input praxicon for recognizing actions via visual-gestural cues, and the output praxicon for producing motor engrams \citep{peigneux2004imaging}. Extensions to AI involve grounding symbolic language in sensorimotor primitives, addressing the symbol grounding problem \citep{harnad1990symbol, pastra2011praxicon, vannuscorps2023effector}.

\subsection{Competing Approaches}
To contextualize our framework, compare it with alternatives. Integrated Information Theory (IIT) posits consciousness as integrated information $\Phi$, but it requires predefined system boundaries and struggles with distributed agency \citep{tononi2004information}. The Free Energy Principle (FEP) frames agents as minimizers of variational free energy, modeling predictive processing but lacking direct measures of structural improbability \citep{friston2010}. Enactive cognition emphasizes autopoiesis and sensorimotor contingencies, yet it often remains qualitative \citep{barandiaran2009defining}. Our EAIMR-praxicon approach complements these by providing computable metrics grounded in assembly and action representation.

\section{Mind Blindness and Intersubjectivity Collapse}
\label{sec:mind_blindness}

Mind blindness, as articulated by Michael Levin, refers to the systematic failure to recognize intelligence in non-human systems, such as plants, animals, and cellular collectives, because they do not conform to human-like symbolic or narrative modes. Levin argues that intelligence is any system’s ability to pursue goals, adapt to perturbations, and solve problems in its medium, evident in regenerating limbs or bioelectric signaling \citep{levin2019computational}. This bias extends beyond ecology to cognitive minorities within humanity, where individuals with aphantasia---the inability to voluntarily create mental images---or anendophasia---the absence of inner monologue---are marginalized despite functional processing capabilities.

\textbf{Mind Blindness Across Domains}

1. \textbf{Ecological Minds (Plants and Animals)} Bias: Intelligence is equated with language and narrative. Blindness: Plants and animals operate through signaling, adaptation, and coordination, but because they do not ``speak,'' their intelligence is dismissed as instinct. Consequence: Systematic marginalization of non-human agency, paving the way for ecological exploitation.

2. \textbf{Cognitive Minorities (Aphantasia, Anendophasia, Atypical Cognition)} Bias: Value is tied to attention performance---vivid self-narration, imagery, persuasion. Blindness: Those whose minds process functionally but without inner imagery or speech must reverse-engineer outputs, making them invisible or ``less capable'' in the attention economy. Aphantasia affects 2-4\% of the population, impairing voluntary visualization, while anendophasia involves silent inner worlds without verbal thought. Consequence: Social and economic exclusion of cognitive styles that do not match platform-optimized norms.

3. \textbf{Religious and Philosophical Systems (Christianity, Islam, Atheism, etc.)} Bias: Each tradition assumes its own narrative or epistemic mode (ritual, revelation, rational critique) is the universal medium of truth. Blindness: Values and intelligences encoded in rival traditions become illegible. What one calls wisdom, another calls heresy or superstition. Consequence: Breakdown of shared meaning across civilizations, escalating into mistrust, conflict, or war.

Unifying Thread: Across all three layers, mind blindness is the refusal to recognize the legitimacy of intelligences that do not mirror our own mode of expression. This blindness is not just a moral failure---it is a structural driver of intersubjectivity collapse, as communication, empathy, and trust fail across ecological, cognitive, and cultural domains. AI futures will magnify this bias---if we only design for intelligences that ``talk like us,'' we will fail to cooperate with (or even recognize) those that do not.

\subsection{Termite--Neuron--Forest Triad: Analogies for Detecting Emergent Intelligence}

In the context of intersubjectivity collapse---where diverse minds fail to communicate or recognize one another due to incompatible cognitive modes---the ``Termite--Neuron--Forest Triad'' serves as a formal analogy to highlight how emergent intelligence arises from swarms of simpler agents. This triad underscores the structural bias of mind blindness, wherein intelligences lacking narrative or imagistic expression are dismissed, exacerbating collapse by excluding valid forms of agency. By applying an assembly-index lens, we can quantify the improbability of these emergent structures, revealing mind-like processes that operate through coordination rather than human-centric symbolism.

1. \textbf{Neurons as ``Just Insects'': The Swarm of Consciousness} Neurons, individually, are simple electrochemical agents following local rules of excitation and inhibition. However, their collective swarming---through synaptic connections, feedback loops, and hierarchical organization---assembles a brain capable of consciousness and goal-directed behavior. The assembly index here is extraordinarily high: the precise mesoscale architecture (e.g., cortical columns) is improbable under mindless diffusion, implying directed processes like plasticity and learning. Mind blindness dismisses this as ``mere biology'' until the emergent whole is recognized, mirroring how we overlook atypical human cognition (e.g., aphantasia) that lacks internal imagery.

2. \textbf{Termites as ``Just Insects'': The Swarm of Architecture} Termites operate as decentralized agents, responding to local pheromones and environmental cues. Yet, their swarming constructs mounds with advanced features: ventilation topology for thermoregulation, chambers for farming fungi, and resilient structures against floods. The mound’s assembly index exceeds that of random soil accumulations (e.g., dunes), reflecting collective intelligence in distributed sensing and adaptation. This parallels neuronal swarming but in an ecological medium, where goals like colony survival manifest without centralized control \citep{turner2000extended, camazine2001self}.

3. \textbf{Trees as ``Just Plants'': The Swarm of Ecosystem Intelligence} Trees, fungi, and associated organisms form a symbiotic swarm, with roots interleaving to share nutrients, signals propagating via mycorrhizal networks, and collective responses to threats like drought or pests. The forest’s assembly index---encompassing trophic cycles, succession dynamics, and long-term resilience---is vastly higher than random vegetation, indicating goal-directed processes at the ecosystem scale. This intelligence, expressed through chemical and bioelectric signaling, evades mind blindness by lacking anthropomorphic traits \citep{levin2019computational}.

Unifying Principle: Across the triad, intelligence emerges from local rules yielding globally improbable assemblies, quantifiable via an effective assembly index adjusted for reuse (motifs like synaptic patterns or root modules), parallelism (concurrent neural firing or fungal growth), and degeneracy (functional alternatives).

Formally, let $A(X)$ denote the raw assembly index of system $X$; the effective index $EA(X) = A(X)/(R\cdot P \cdot D)$, where $R$, $P$, and $D$ are reuse, parallelism, and degeneracy factors, signals agency when exceeding thresholds derived from null models. This bypasses narrativity bias, aligning with Levin’s emphasis on goal pursuit in non-brain systems. In intersubjectivity collapse, unrecognized triadic minds accelerate fragmentation.

To extend this triad, we incorporate plant cognition research, where mycorrhizal networks exhibit memory and decision-making, further challenging mind blindness \citep{calvo2020plants}. Collective intelligence studies show how swarms scale agency \citep{camazine2001self}, while network neuroscience reveals similar patterns in brain connectivity \citep{bassett2017network}. Indigenous knowledge systems offer non-anthropocentric views of agency in nature \citep{kimmerer2013braiding}, enriching the framework with extended cognition perspectives \citep{clark2010supersizing}.

\section{Formal Definitions}
\label{sec:formal}

\subsection{Prerequisites: Probabilistic Models and Generative Distributions}
Understanding EAIMR requires familiarity with generative models. A null model $P_0(X)$ represents mindless processes, such as stochastic diffusion or self-organization without goals, while a mind-like model $P_1(X)$ incorporates agency, such as optimization toward functional setpoints. These distributions allow computation of likelihoods and surprises, foundational to information-theoretic analyses \citep{cover2006elements}.

\subsection{EAIMR Definition}
For a system $X$ described over an alphabet $\Sigma$, EAIMR is:

\[ \text{EAIMR}(X) = \frac{A(X)}{R(X) \cdot P(X) \cdot D(X)} \cdot C(X), \]

where: - $A(X)$: Raw Assembly Index, the minimal steps in an assembly grammar \citep{walker2013algorithmic}. - $R(X)$: Reuse factor, quantifying modular repetition. - $P(X)$: Parallelism factor, accounting for concurrent processes. - $D(X)$: Degeneracy factor, reflecting alternative configurations achieving similar functions. - $C(X)$: Coherence factor, amplifying goal-directed integration. High values suggest agency, as the structure is improbable under $P_0(X)$ but plausible under $P_1(X)$.

\subsection{Praxicon Definition}
The praxicon is formalized as a bipartite graph $G_P = (V_A, V_S, E)$, with action vertices $V_A$ (e.g., motor programs), semantic vertices $V_S$ (e.g., conceptual labels), and edges $E$ encoding mappings \citep{panunzi2018imagact}. This structure facilitates bidirectional translation between perception and action, essential for embodied cognition \citep{barsalou2008grounded}.

\section{Mathematical Integration}
\label{sec:integration}

\subsection{Information-Theoretic Foundations}
Self-information under the null: $I_0(X) = -\log P_0(X)$. The likelihood ratio $\Lambda(X) = P_1(X)/P_0(X)$ quantifies preference for agentic models. For praxicon integration, action entropy reduction is:

$\Delta H_P = H(P_0) - H_{\text{obs}}(P)$,

updating coherence as $C(X) \approx 1 + \alpha \Delta H_P / H(P_0)$ \citep{shannon1948mathematical, jaynes2003probability}. Kolmogorov complexity approximates $A(X)$ \citep{kolmogorov1965three}.

\subsection{Graph-Theoretic Formalism}
The praxicon graph $G_P$ allows coherence via spectral properties, e.g., eigenvalue stability for robustness. EAIMR identifies low-entropy subgraphs indicative of agency \citep{newell1990unified}.

\subsection{Category-Theoretic Interpretation}
Category theory provides a formal language for describing structures and transformations in abstract terms, making it particularly suitable for modeling cognitive processes that involve mappings between different domains \citep{healy2022, phillips2021}. A category consists of objects (e.g., perceptual states, motor actions) and morphisms (transformations or relations between objects) that satisfy associativity and identity axioms. Functors map between categories while preserving structure, and natural transformations provide morphisms between functors. In our framework, the praxicon can be viewed as a functor $F : C_{\text{perc}} \to C_{\text{motor}}$, where $C_{\text{perc}}$ is the category of perceptual inputs (objects as sensory data, morphisms as perceptual transformations) and $C_{\text{motor}}$ is the category of motor outputs (objects as action states, morphisms as motor adjustments), mediated through a semantic category $C_{\text{sem}}$. This functorial representation captures how perceptual inputs are systematically transformed into motor actions via semantic grounding. EAIMR, in this context, quantifies the ``improbability'' or cost of natural transformations $\eta : F \implies G$, where $G$ represents alternative mappings (e.g., under null dynamics). High EAIMR corresponds to scenarios where the observed natural transformation (the agentic mapping) is highly specific and unlikely under random perturbations, reflecting the coordinated agency. For instance, in neural systems, category-theoretic models can describe how cognitive development involves composing functors for increasingly complex action representations \citep{ehresmann1980}. This interpretation aligns with structuralist views in cognitive science, emphasizing relational structures over material substrates \citep{ehresmann2021}.

\subsection{Sheaf-Theoretic Interpretation}
Sheaf theory extends category theory by addressing how local data on a topological space can be consistently glued into global structures, offering a powerful tool for modeling distributed systems where agency emerges from multi-scale interactions \citep{spivak2024}. A sheaf $F$ on a space $X$ assigns to each open set $U \subseteq X$ a set $F(U)$ (local data) with restriction maps that satisfy gluing conditions: local sections that agree on overlaps can be uniquely combined into a global section. In biological and ecological contexts, sheaf theory models uncertainty and resilience in complex networks \citep{stenzel2024}. For our framework, define a sheaf $F$ on the topological space representing the system’s components (e.g., neurons in a circuit or species in an ecosystem), where sections over open sets encode local praxicon fragments---action representations valid in subsystems. The gluing axiom ensures that local actions cohere into global behaviors, such as coordinated foraging in ant colonies or nutrient cycling in forests. EAIMR quantifies obstructions to trivial gluing, often captured by sheaf cohomology groups $H^k(X, F)$. Non-zero cohomology indicates ``twists'' or global inconsistencies under null models, corresponding to high improbability and thus agency. For example, in ecology, forest nutrient networks form a sheaf where local flows (sections) glue into global cycles; high $H^1$ (non-exact cycles) elevates EAIMR, detecting emergent intelligence \citep{levin2019computational}. In biology, bioelectric signaling can be sheafified, with cohomology measuring scale-free cognition \citep{levin2019computational, spivak2024}. This approach bridges local computations to global agency, providing a mathematical basis for resilience under perturbations.

\section{Worked Examples}
\label{sec:examples}

\subsection{Neural Tissue}
In neural systems, high $A(X)$ arises from precise synaptic architectures, with praxicon engrams in cortical columns reducing entropy. Simulations of mirror neurons show effector-specific mappings \citep{gallese1998mirror}. Step-by-step: Compute $A(X)$ via connection counts, adjust discounts, amplify with coherence from goal pursuit. Entropy curve: Sharp reduction in action space.

\subsection{Termite Mound}
Moderate assembly in tunnels, high coherence in thermoregulation. Collective praxicon via pheromone-guided motifs \citep{turner2000extended}. Graph: Nodes for worker actions linked to colony semantics. Data from self-organization models \citep{camazine2001self}.

\subsection{Forest Ecosystem}
Extensive macro-assembly in symbiotic networks, praxicon-like routing. High degeneracy in species roles, coherence over timescales \citep{levin2019computational}. Lattice diagram: Interconnected trophic levels.

\subsection{Negative Case: Crystal Growth}
Regular lattice growth yields low $C(X)$, high predictability under $P_0$, resulting in low EAIMR \citep{hazen2007functional}.

\section{Theoretical Extensions}
\label{sec:theoretical_extensions}

To enhance the framework, we introduce multi-scale temporal dynamics with a temporal coherence component $C_t(X)$, measuring agency across timescales from neural firing to forest succession, distinguishing long-term agency from patterns \citep{levin2019computational}.

Causal intervention frameworks define an agency resilience index, quantifying recovery from targeted disruptions versus noise, differentiating true agents from homeostatic systems \citep{barandiaran2009defining}.

Compositional agency extends category theory for nested intelligences, where agents form collectives with irreducible properties \citep{ehresmann2021}.

\section{Empirical and Methodological Improvements}
\label{sec:empirical}

Quantitative validation studies are essential, including neural comparisons of healthy vs. lesioned tissue, ecological manipulations of forests/termite colonies, AI benchmarks, and cross-cultural ritual analyses \citep{pastra2011praxicon}.

Computational implementation requires open-source tools for EAIMR computation, with standardized data formats, Kolmogorov approximation algorithms, and praxicon graph extraction methods \citep{kolmogorov1965three}.

Statistical framework includes confidence intervals on EAIMR, multiple comparison corrections, power analysis for agency detection, and handling uncertainty \citep{jaynes2003probability}.

\section{Conceptual Deepening}
\label{sec:conceptual}

The phenomenological bridge addresses consciousness: high EAIMR systems may exhibit phenomenal experience, with testable predictions for subjective states \citep{tononi2004information}.

Evolutionary dynamics predict higher EAIMR systems show distinct trajectories, testable in simulations or comparative studies \citep{camazine2001self}.

Information-theoretic refinements connect EAIMR to integrated information, non-equilibrium thermodynamics, and predictive processing \citep{friston2010}.

\section{Practical Applications}
\label{sec:applications}

\subsection{Neuroscience}
Detects apraxia and atypical cognition \citep{peigneux2004imaging, zeman2015}. Clinical applications include diagnostics for locked-in syndrome, rehabilitation monitoring, and therapeutic community agency assessment \citep{levin2022technological}.

\subsection{AI and Robotics}
Grounds LLMs; detects agency \citep{pastra2011praxicon, anderson2007how}. AI alignment protocols monitor emergent agency during training, detect deceptive alignment, and measure value alignment via praxicon analysis \citep{yudkowsky2008artificial}.

\subsection{Ecological and Collective Minds}
Identifies swarm intelligence \citep{camazine2001self, simondon2017mode}. Policy frameworks enable legal recognition of non-human agency, environmental protection based on ecosystem metrics, and AI regulation \citep{bostrom2014superintelligence}.

\subsection{Cross-Cultural Cognition}
Rituals as praxicons \citep{ortega1961meditations}. Anthropocentric bias mitigation involves cross-validation with non-human measures, indigenous knowledge, and animal studies \citep{kimmerer2013braiding}.

\section{Limits and Future Research}
\label{sec:limits}

Challenges: Approximating complexity, validating sheaves. Computational tractability requires polynomial approximations and hierarchical methods for large systems \citep{kolmogorov1965three}.

Future: Bioelectric models, AI simulations \citep{levin2019computational}. Connections to plant cognition \citep{calvo2020plants}, collective intelligence \citep{camazine2001self}, network neuroscience \citep{bassett2017network}, indigenous knowledge \citep{kimmerer2013braiding}, and extended cognition \citep{clark2010supersizing}.

\section{Ethical and Philosophical Implications}
\label{sec:ethics}

Mitigates mind blindness, informs AI ethics, and aligns with process philosophy \citep{dennett1987intentional, whitehead1929process, deleuze1987thousand}.

\section{Decision Rule}
\label{sec:decision}

$\text{EAIMR}(X) \ge \tau$ classifies agency, using graph entropy proxies.

\section{Conclusion}
\label{sec:conclusion}

This framework detects agency, averting intersubjectivity collapse by recognizing diverse minds.

\appendix

\section{Mathematical Appendix: EAIMR, Information, and Improbable Order}

A.1 Preliminaries and Notation

Let $X$ denote a structure/system (e.g., termite mound, forest network, neural microcircuit) described over an alphabet $\Sigma$. Let $P_0$ be a null (mindless) generative model for $X$ (e.g., stochastic physics + simple self-organization) with distribution $P_0(X)$. Let $P_1$ be a mind-like generative model for $X$ (agentic/goal-directed) with distribution $P_1(X)$. Raw Assembly Index $A(X)$: minimum step count in a prescribed assembly grammar to build $X$ from primitives.

EAIMR (from your definition):

\[ \text{EAIMR}(X) = \frac{A(X)}{R(X) \cdot P(X) \cdot D(X) \cdot C(X)} \]

A.2 Connection to Kolmogorov Complexity

Let $K(X)$ be prefix Kolmogorov complexity (shortest universal description of $X$). In general, lower $K(X)$ higher compressibility (simple regularities). Many mind-built artifacts have mixed profiles: globally compressible high-level plan + low-level details that are costly to specify. A useful proxy is two-part MDL (minimum description length):

\[ L(X) \approx L(\text{model}) + L(\text{residuals} \mid \text{model}) \]

Heuristic link (non-rigorous but actionable):

\[ \text{EAIMR}(X) \propto \frac{K_{\text{plan}}(X)}{R \cdot P \cdot D} \cdot C \]

with $K_{\text{plan}}(X) \lesssim K(X)$.

Kolmogorov complexity $K(X)$ is the length of the shortest program $p$ in a fixed universal Turing machine $U$ such that $U(p) = X$. For EAIMR, we approximate $K_{\text{plan}}(X)$ as the compressible plan component, separating it from residuals using two-part coding, which aligns with practical compression algorithms like LZW or gzip for estimation.

A.3 Statistical Surprise and Likelihood Ratios

Define the self-information under the null:

\[ I_0(X) \equiv -\log P_0(X) \]

\[ \Lambda(X) \equiv \frac{P_1(X)}{P_0(X)} = \exp \left( \log P_1(X) - \log P_0(X) \right). \]

Key idea: High EAIMR should correlate with large $\Lambda(X)$ and a high $I_0(X)$ (i.e., $X$ is very unlikely under mindless dynamics but plausible under mind-like dynamics). A practical surrogate score you can compute or estimate:

\[ S(X) \equiv I_0(X) + \lambda \phi(X) \]

\[ \text{EAIMR}(X) = g \left( S(X) \right), \quad g \uparrow \]

The likelihood ratio $\Lambda(X)$ is the Bayes factor for mind-like vs. null models, with $\log \Lambda(X)$ as evidence weight. For estimation, use Monte Carlo methods to approximate $P_0(X)$ and $P_1(X)$ from generative simulations, such as Markov chain Monte Carlo for high-dimensional spaces.

A.4 Entropy Reduction, Work, and Directed Assembly

Let $M$ be macrostate features (e.g., ventilation topology, nutrient routing, neural coding). Define macro-entropy under $P_0$ and observed for $X$. Directed assembly manifests as entropy reduction:

\[ \Delta H \equiv H(M) - H_{\text{obs}}(M) > 0 \]

A convenient coherence multiplier:

\[ C(X) \approx 1 + \alpha \frac{\Delta H}{H(M)} \]

Entropy $H(M) = -\sum p_i \log p_i$ over macrostates, where $p_i$ is probability under $P_0$. Observed $H_{\text{obs}}(M)$ is computed from empirical distributions, with $\Delta H$ quantifying work against thermodynamic disorder, linking to non-equilibrium physics \citep{friston2010}.

A.5 Reuse, Parallelism, Degeneracy as Information Discounts

Treat these as code-length discounts relative to the raw assembly plan. Reuse: library motifs shorten description. If motif dictionary has description amortized across instances, the effective bits per instance drop. Model as a divisor on $A(X)$. Parallelism: independent sub-assemblies compress as a product code; wall-clock and effective complexity per unit time fall. Model as a divisor on $A(X)$. Degeneracy: if many alternatives achieve the same function, the description of ``a solution'' is shorter than ``this specific solution.'' Function-first coding reduces effective complexity; model as a divisor. This aligns your original formula with MDL-style accounting: EAIMR rises when, despite these discounts, the remaining coordinated specification is still large and functionally tight.

For reuse, $R(X) = 1 + \log(\text{number of motif instances} / \text{motif types})$, amortizing costs. Parallelism $P(X) = \text{total operations} / \text{critical path length}$, capturing concurrency. Degeneracy $D(X) = \log(\text{number of functional equivalents})$, spreading probability mass.

A.6 Coherence via Control and Causal Graphs

Let $G(X)$ be a causal graph extracted from $X$’s dynamics. Define a control coherence term:

\[ \phi(X) = \beta_1 \text{Stab}(G) + \beta_2 \text{GoalAcc}(G) + \beta_3 \text{Robust}(G) \]

GoalAcc($G$): accuracy tracking a specified functional setpoint (e.g., temperature band in termite mounds). Robust($G$): performance under noise/lesions (graceful degradation). Then use $\phi(X)$ inside (Section 3) or directly as your multiplier.

Stab($G$) can be graph Laplacian eigenvalues for stability; GoalAcc($G$) as mean squared error from setpoint; Robust($G$) as fraction of performance retained under simulated lesions, using graph pruning algorithms.

A.7 Practical Estimation (finite data)

Because $K(X)$ is uncomputable, rely on approximations: Compression proxies: on suitable encodings (graphs, fields, symbol streams). Surprisal under null: fit $P_0$ (e.g., generative physics sims, random graphs) and estimate by Monte Carlo or density models. Function and coherence: quantify with control/evaluation tasks (tracking error, energy efficiency, resilience curves). Discounts: measure from libraries (motif counts), concurrency (critical path vs. total ops), and function-equivalence classes (number of viable variants).

For compression, use BZ2 or LZMA ratios as $K(X)$ lower bounds. Monte Carlo for $P_0$ uses $10^4-10^6$ samples for convergence. Resilience curves fit to exponential decay models for Robust($G$).

A.8 Worked Intuitions (triad examples)

Neural tissue: very high $A(X)$ (precise mesoscale organization), strong $C(X)$ (robust goal pursuit), discounts present (columns, parallelism), yet EAIMR stays high mind-like. Termite mound: moderate $A(X)$ vs. dunes, high $C(X)$ (thermoregulation), sizable discounts (motifs, parallel labor, many workable designs) but coordinated control keeps EAIMR high. Forest: extremely high macro $A(X)$ vs. random vegetation, high robustness/goal achievement (nutrient routing, succession), massive reuse/degeneracy, yet coherence over centuries pushes EAIMR high.

A.9 Decision Rule (operational)

Define thresholds from a reference set of clearly mindless assemblies:

$\text{EAIMR}(X) \ge \tau \implies$ classify as mind-like (candidate)

--- Takeaway

Information-theoretic view: EAIMR is a mind detector for improbably coherent order---high surprisal under a null, discounted for shortcuts, and boosted by functional integration. It sidesteps narrativity bias: you do not need language or imagery to register as intelligent; termite colonies, forests, and atypical human cognition can score high if they assemble and maintain low-entropy, goal-directed structure against default dynamics.

B Derivations and Pseudocode

B.1 EAIMR Derivation

From surprisal: $S(X) = I_0(X) + \lambda \phi(X)$, EAIMR as monotonic $g(S(X))$.

B.2 Pseudocode

\begin{verbatim}
def compute_eaimr(X, P0, P1, G_P):
    A = assembly_index(X)
    R = reuse_factor(X)
    P = parallelism(X)
    D = degeneracy(X)
    C = coherence(G_P)
    return (A / (R * P * D)) * C
\end{verbatim}

\bibliographystyle{plainnat}
\bibliography{references}

\end{document}
where: - $A(X)$: Raw Assembly Index, the minimal steps in an assembly grammar \citep{walker2013algorithmic}. - $R(X)$: Reuse factor, quantifying modular repetition. - $P(X)$: Parallelism factor, accounting for concurrent processes. - $D(X)$: Degeneracy factor, reflecting alternative configurations achieving similar functions. - $C(X)$: Coherence factor, amplifying goal-directed integration. High values suggest agency, as the structure is improbable under $P_0(X)$ but plausible under $P_1(X)$.

\subsection{Praxicon Definition}
The praxicon is formalized as a bipartite graph $G_P = (V_A, V_S, E)$, with action vertices $V_A$ (e.g., motor programs), semantic vertices $V_S$ (e.g., conceptual labels), and edges $E$ encoding mappings \citep{panunzi2018imagact}. This structure facilitates bidirectional translation between perception and action, essential for embodied cognition \citep{barsalou2008grounded}.

\section{Mathematical Integration}
\label{sec:integration}

\subsection{Information-Theoretic Foundations}
Self-information under the null: $I_0(X) = -\log P_0(X)$. The likelihood ratio $\Lambda(X) = P_1(X)/P_0(X)$ quantifies preference for agentic models. For praxicon integration, action entropy reduction is:

$\Delta H_P = H(P_0) - H_{\text{obs}}(P)$,

updating coherence as $C(X) \approx 1 + \alpha \Delta H_P / H(P_0)$ \citep{shannon1948mathematical, jaynes2003probability}. Kolmogorov complexity approximates $A(X)$ \citep{kolmogorov1965three}.

\subsection{Graph-Theoretic Formalism}
The praxicon graph $G_P$ allows coherence via spectral properties, e.g., eigenvalue stability for robustness. EAIMR identifies low-entropy subgraphs indicative of agency \citep{newell1990unified}.

\subsection{Category-Theoretic Interpretation}
Category theory provides a formal language for describing structures and transformations in abstract terms, making it particularly suitable for modeling cognitive processes that involve mappings between different domains \citep{healy2022, phillips2021}. A category consists of objects (e.g., perceptual states, motor actions) and morphisms (transformations or relations between objects) that satisfy associativity and identity axioms. Functors map between categories while preserving structure, and natural transformations provide morphisms between functors. In our framework, the praxicon can be viewed as a functor $F : C_{\text{perc}} \to C_{\text{motor}}$, where $C_{\text{perc}}$ is the category of perceptual inputs (objects as sensory data, morphisms as perceptual transformations) and $C_{\text{motor}}$ is the category of motor outputs (objects as action states, morphisms as motor adjustments), mediated through a semantic category $C_{\text{sem}}$. This functorial representation captures how perceptual inputs are systematically transformed into motor actions via semantic grounding. EAIMR, in this context, quantifies the ``improbability'' or cost of natural transformations $\eta : F \implies G$, where $G$ represents alternative mappings (e.g., under null dynamics). High EAIMR corresponds to scenarios where the observed natural transformation (the agentic mapping) is highly specific and unlikely under random perturbations, reflecting the coordinated agency. For instance, in neural systems, category-theoretic models can describe how cognitive development involves composing functors for increasingly complex action representations \citep{ehresmann1980}. This interpretation aligns with structuralist views in cognitive science, emphasizing relational structures over material substrates \citep{ehresmann2021}.

\subsection{Sheaf-Theoretic Interpretation}
Sheaf theory extends category theory by addressing how local data on a topological space can be consistently glued into global structures, offering a powerful tool for modeling distributed systems where agency emerges from multi-scale interactions \citep{spivak2024}. A sheaf $F$ on a space $X$ assigns to each open set $U \subseteq X$ a set $F(U)$ (local data) with restriction maps that satisfy gluing conditions: local sections that agree on overlaps can be uniquely combined into a global section. In biological and ecological contexts, sheaf theory models uncertainty and resilience in complex networks \citep{stenzel2024}. For our framework, define a sheaf $F$ on the topological space representing the system’s components (e.g., neurons in a circuit or species in an ecosystem), where sections over open sets encode local praxicon fragments---action representations valid in subsystems. The gluing axiom ensures that local actions cohere into global behaviors, such as coordinated foraging in ant colonies or nutrient cycling in forests. EAIMR quantifies obstructions to trivial gluing, often captured by sheaf cohomology groups $H^k(X, F)$. Non-zero cohomology indicates ``twists'' or global inconsistencies under null models, corresponding to high improbability and thus agency. For example, in ecology, forest nutrient networks form a sheaf where local flows (sections) glue into global cycles; high $H^1$ (non-exact cycles) elevates EAIMR, detecting emergent intelligence \citep{levin2019computational}. In biology, bioelectric signaling can be sheafified, with cohomology measuring scale-free cognition \citep{levin2019computational, spivak2024}. This approach bridges local computations to global agency, providing a mathematical basis for resilience under perturbations.

\section{Worked Examples}
\label{sec:examples}

\subsection{Neural Tissue}
In neural systems, high $A(X)$ arises from precise synaptic architectures, with praxicon engrams in cortical columns reducing entropy. Simulations of mirror neurons show effector-specific mappings \citep{gallese1998mirror}. Step-by-step: Compute $A(X)$ via connection counts, adjust discounts, amplify with coherence from goal pursuit. Entropy curve: Sharp reduction in action space.

\subsection{Termite Mound}
Moderate assembly in tunnels, high coherence in thermoregulation. Collective praxicon via pheromone-guided motifs \citep{turner2000extended}. Graph: Nodes for worker actions linked to colony semantics. Data from self-organization models \citep{camazine2001self}.

\subsection{Forest Ecosystem}
Extensive macro-assembly in symbiotic networks, praxicon-like routing. High degeneracy in species roles, coherence over timescales \citep{levin2019computational}. Lattice diagram: Interconnected trophic levels.

\subsection{Negative Case: Crystal Growth}
Regular lattice growth yields low $C(X)$, high predictability under $P_0$, resulting in low EAIMR \citep{hazen2007functional}.

\section{Theoretical Extensions}
\label{sec:theoretical_extensions}

To enhance the framework, we introduce multi-scale temporal dynamics with a temporal coherence component $C_t(X)$, measuring agency across timescales from neural firing to forest succession, distinguishing long-term agency from patterns \citep{levin2019computational}.

Causal intervention frameworks define an agency resilience index, quantifying recovery from targeted disruptions versus noise, differentiating true agents from homeostatic systems \citep{barandiaran2009defining}.

Compositional agency extends category theory for nested intelligences, where agents form collectives with irreducible properties \citep{ehresmann2021}.

\section{Empirical and Methodological Improvements}
\label{sec:empirical}

Quantitative validation studies are essential, including neural comparisons of healthy vs. lesioned tissue, ecological manipulations of forests/termite colonies, AI benchmarks, and cross-cultural ritual analyses \citep{pastra2011praxicon}.

Computational implementation requires open-source tools for EAIMR computation, with standardized data formats, Kolmogorov approximation algorithms, and praxicon graph extraction methods \citep{kolmogorov1965three}.

Statistical framework includes confidence intervals on EAIMR, multiple comparison corrections, power analysis for agency detection, and handling uncertainty \citep{jaynes2003probability}.

\section{Conceptual Deepening}
\label{sec:conceptual}

The phenomenological bridge addresses consciousness: high EAIMR systems may exhibit phenomenal experience, with testable predictions for subjective states \citep{tononi2004information}.

Evolutionary dynamics predict higher EAIMR systems show distinct trajectories, testable in simulations or comparative studies \citep{camazine2001self}.

Information-theoretic refinements connect EAIMR to integrated information, non-equilibrium thermodynamics, and predictive processing \citep{friston2010}.

\section{Practical Applications}
\label{sec:applications}

\subsection{Neuroscience}
Detects apraxia and atypical cognition \citep{peigneux2004imaging, zeman2015}. Clinical applications include diagnostics for locked-in syndrome, rehabilitation monitoring, and therapeutic community agency assessment \citep{levin2022technological}.

\subsection{AI and Robotics}
Grounds LLMs; detects agency \citep{pastra2011praxicon, anderson2007how}. AI alignment protocols monitor emergent agency during training, detect deceptive alignment, and measure value alignment via praxicon analysis \citep{yudkowsky2008artificial}.

\subsection{Ecological and Collective Minds}
Identifies swarm intelligence \citep{camazine2001self, simondon2017mode}. Policy frameworks enable legal recognition of non-human agency, environmental protection based on ecosystem metrics, and AI regulation \citep{bostrom2014superintelligence}.

\subsection{Cross-Cultural Cognition}
Rituals as praxicons \citep{ortega1961meditations}. Anthropocentric bias mitigation involves cross-validation with non-human measures, indigenous knowledge, and animal studies \citep{kimmerer2013braiding}.

\section{Limits and Future Research}
\label{sec:limits}

Challenges: Approximating complexity, validating sheaves. Computational tractability requires polynomial approximations and hierarchical methods for large systems \citep{kolmogorov1965three}.

Future: Bioelectric models, AI simulations \citep{levin2019computational}. Connections to plant cognition \citep{calvo2020plants}, collective intelligence \citep{camazine2001self}, network neuroscience \citep{bassett2017network}, indigenous knowledge \citep{kimmerer2013braiding}, and extended cognition \citep{clark2010supersizing}.

\section{Ethical and Philosophical Implications}
\label{sec:ethics}

Mitigates mind blindness, informs AI ethics, and aligns with process philosophy \citep{dennett1987intentional, whitehead1929process, deleuze1987thousand}.

\section{Decision Rule}
\label{sec:decision}

$\text{EAIMR}(X) \ge \tau$ classifies agency, using graph entropy proxies.

\section{Conclusion}
\label{sec:conclusion}

This framework detects agency, averting intersubjectivity collapse by recognizing diverse minds.

\appendix

\section{Mathematical Appendix: EAIMR, Information, and Improbable Order}

A.1 Preliminaries and Notation

Let $X$ denote a structure/system (e.g., termite mound, forest network, neural microcircuit) described over an alphabet $\Sigma$. Let $P_0$ be a null (mindless) generative model for $X$ (e.g., stochastic physics + simple self-organization) with distribution $P_0(X)$. Let $P_1$ be a mind-like generative model for $X$ (agentic/goal-directed) with distribution $P_1(X)$. Raw Assembly Index $A(X)$: minimum step count in a prescribed assembly grammar to build $X$ from primitives.

EAIMR (from your definition):

\[ \text{EAIMR}(X) = \frac{A(X)}{R(X) \cdot P(X) \cdot D(X) \cdot C(X)} \]

A.2 Connection to Kolmogorov Complexity

Let $K(X)$ be prefix Kolmogorov complexity (shortest universal description of $X$). In general, lower $K(X)$ higher compressibility (simple regularities). Many mind-built artifacts have mixed profiles: globally compressible high-level plan + low-level details that are costly to specify. A useful proxy is two-part MDL (minimum description length):

\[ L(X) \approx L(\text{model}) + L(\text{residuals} \mid \text{model}) \]

Heuristic link (non-rigorous but actionable):

\[ \text{EAIMR}(X) \propto \frac{K_{\text{plan}}(X)}{R \cdot P \cdot D} \cdot C \]

with $K_{\text{plan}}(X) \lesssim K(X)$.

Kolmogorov complexity $K(X)$ is the length of the shortest program $p$ in a fixed universal Turing machine $U$ such that $U(p) = X$. For EAIMR, we approximate $K_{\text{plan}}(X)$ as the compressible plan component, separating it from residuals using two-part coding, which aligns with practical compression algorithms like LZW or gzip for estimation.

A.3 Statistical Surprise and Likelihood Ratios

Define the self-information under the null:

\[ I_0(X) \equiv -\log P_0(X) \]

\[ \Lambda(X) \equiv \frac{P_1(X)}{P_0(X)} = \exp \left( \log P_1(X) - \log P_0(X) \right). \]

Key idea: High EAIMR should correlate with large $\Lambda(X)$ and a high $I_0(X)$ (i.e., $X$ is very unlikely under mindless dynamics but plausible under mind-like dynamics). A practical surrogate score you can compute or estimate:

\[ S(X) \equiv I_0(X) + \lambda \phi(X) \]

\[ \text{EAIMR}(X) = g \left( S(X) \right), \quad g \uparrow \]

The likelihood ratio $\Lambda(X)$ is the Bayes factor for mind-like vs. null models, with $\log \Lambda(X)$ as evidence weight. For estimation, use Monte Carlo methods to approximate $P_0(X)$ and $P_1(X)$ from generative simulations, such as Markov chain Monte Carlo for high-dimensional spaces.

A.4 Entropy Reduction, Work, and Directed Assembly

Let $M$ be macrostate features (e.g., ventilation topology, nutrient routing, neural coding). Define macro-entropy under $P_0$ and observed for $X$. Directed assembly manifests as entropy reduction:

\[ \Delta H \equiv H(M) - H_{\text{obs}}(M) > 0 \]

A convenient coherence multiplier:

\[ C(X) \approx 1 + \alpha \frac{\Delta H}{H(M)} \]

Entropy $H(M) = -\sum p_i \log p_i$ over macrostates, where $p_i$ is probability under $P_0$. Observed $H_{\text{obs}}(M)$ is computed from empirical distributions, with $\Delta H$ quantifying work against thermodynamic disorder, linking to non-equilibrium physics \citep{friston2010}.

A.5 Reuse, Parallelism, Degeneracy as Information Discounts

Treat these as code-length discounts relative to the raw assembly plan. Reuse: library motifs shorten description. If motif dictionary has description amortized across instances, the effective bits per instance drop. Model as a divisor on $A(X)$. Parallelism: independent sub-assemblies compress as a product code; wall-clock and effective complexity per unit time fall. Model as a divisor on $A(X)$. Degeneracy: if many alternatives achieve the same function, the description of ``a solution'' is shorter than ``this specific solution.'' Function-first coding reduces effective complexity; model as a divisor. This aligns your original formula with MDL-style accounting: EAIMR rises when, despite these discounts, the remaining coordinated specification is still large and functionally tight.

For reuse, $R(X) = 1 + \log(\text{number of motif instances} / \text{motif types})$, amortizing costs. Parallelism $P(X) = \text{total operations} / \text{critical path length}$, capturing concurrency. Degeneracy $D(X) = \log(\text{number of functional equivalents})$, spreading probability mass.

A.6 Coherence via Control and Causal Graphs

Let $G(X)$ be a causal graph extracted from $X$’s dynamics. Define a control coherence term:

\[ \phi(X) = \beta_1 \text{Stab}(G) + \beta_2 \text{GoalAcc}(G) + \beta_3 \text{Robust}(G) \]

GoalAcc($G$): accuracy tracking a specified functional setpoint (e.g., temperature band in termite mounds). Robust($G$): performance under noise/lesions (graceful degradation). Then use $\phi(X)$ inside (Section 3) or directly as your multiplier.

Stab($G$) can be graph Laplacian eigenvalues for stability; GoalAcc($G$) as mean squared error from setpoint; Robust($G$) as fraction of performance retained under simulated lesions, using graph pruning algorithms.

A.7 Practical Estimation (finite data)

Because $K(X)$ is uncomputable, rely on approximations: Compression proxies: on suitable encodings (graphs, fields, symbol streams). Surprisal under null: fit $P_0$ (e.g., generative physics sims, random graphs) and estimate by Monte Carlo or density models. Function and coherence: quantify with control/evaluation tasks (tracking error, energy efficiency, resilience curves). Discounts: measure from libraries (motif counts), concurrency (critical path vs. total ops), and function-equivalence classes (number of viable variants).

For compression, use BZ2 or LZMA ratios as $K(X)$ lower bounds. Monte Carlo for $P_0$ uses $10^4-10^6$ samples for convergence. Resilience curves fit to exponential decay models for Robust($G$).

A.8 Worked Intuitions (triad examples)

Neural tissue: very high $A(X)$ (precise mesoscale organization), strong $C(X)$ (robust goal pursuit), discounts present (columns, parallelism), yet EAIMR stays high mind-like. Termite mound: moderate $A(X)$ vs. dunes, high $C(X)$ (thermoregulation), sizable discounts (motifs, parallel labor, many workable designs) but coordinated control keeps EAIMR high. Forest: extremely high macro $A(X)$ vs. random vegetation, high robustness/goal achievement (nutrient routing, succession), massive reuse/degeneracy, yet coherence over centuries pushes EAIMR high.

A.9 Decision Rule (operational)

Define thresholds from a reference set of clearly mindless assemblies:

$\text{EAIMR}(X) \ge \tau \implies$ classify as mind-like (candidate)

--- Takeaway

Information-theoretic view: EAIMR is a mind detector for improbably coherent order---high surprisal under a null, discounted for shortcuts, and boosted by functional integration. It sidesteps narrativity bias: you do not need language or imagery to register as intelligent; termite colonies, forests, and atypical human cognition can score high if they assemble and maintain low-entropy, goal-directed structure against default dynamics.

B Derivations and Pseudocode

B.1 EAIMR Derivation

From surprisal: $S(X) = I_0(X) + \lambda \phi(X)$, EAIMR as monotonic $g(S(X))$.

B.2 Pseudocode

\begin{verbatim}
def compute_eaimr(X, P0, P1, G_P):
    A = assembly_index(X)
    R = reuse_factor(X)
    P = parallelism(X)
    D = degeneracy(X)
    C = coherence(G_P)
    return (A / (R * P * D)) * C
\end{verbatim}

\bibliographystyle{plainnat}
\bibliography{references}

\end{document}
